\section{Explicit finite checks and the limits of computation}
\label{app:finite}
\subsection{A two-state posterior calculation}
The smallest dynamic instance has $d=1$, $\gamma=4/5$, and
$\epsilon=1/4$. Put $v=\tanh u$ and select $v=1/2$, which corresponds
to the legal command $u=\operatorname{artanh}(1/2)<1$. Starting from the
uniform preparation, the pre-report belief is $q=(1/2,1/2)$, and
\[
 \Pp(Y=+1)=\frac{17}{32},\qquad
 \Pp(Y=-1)=\frac{15}{32}.
\]
The posterior probabilities of state $1$ are respectively
\[
 p_1^+=\frac9{17},\qquad p_1^-=\frac7{15}.
\]
Thus the same current command produces two different posteriors. Discarding
the negative report or assigning its conditional law mass one would alter
the acquisition law used in the proof. With the fixed terminal probe the
corresponding success probabilities are
\[
 \frac12+\frac14\frac9{17}=\frac{43}{68},\qquad
 \frac12+\frac14\frac7{15}=\frac{37}{60}.
\]
The mean posterior agrees with the pre-report probability:
\[
 \frac{17}{32}\frac9{17}+\frac{15}{32}\frac7{15}=\frac12.
\]
This is the posterior martingale property for the report step; the refresh
step has its own affine prediction and is not being omitted.

The derivative with respect to $v$, rather than $u$, at this point for the
positive report is
\[
 \frac{q_0q_1\epsilon}{D_+^2}
 =\frac{16}{289}.
\]
Multiplying by $dv/du=1-v^2=3/4$ gives $12/289$, exactly
\eqref{eq:refresh-jacobian} in one dimension. These distinct derivatives
are recorded because replacing a physical command by a nonlinear
coordinate without its Jacobian would change the preparation density.

\subsection{A direct resource calculation}
Suppose a report converter has two persistent states (for example, it
retains the previous binary report), the target predictor has $M$ labels,
and a feedback controller has $D$ labels. The implemented composite uses
at most $2MD$ triples. On the next step it reads the old triple, the current
input, and fresh independent randomness, and writes the next triple.
The source history may be longer, but it is not consulted by this
implementation. Any reduction below $2MD$ needs a proved state
identification; it is not obtained by calling the converter a sufficient
statistic.

Conversely, if a converter is memoryless and deterministic, its state set
has one element. In that special case it need not enlarge the target
label count. This is the correct context for a nonincrease assertion.
It does not imply the false general statement about exact predictive
memory refuted in Example~\ref{ex:bit-gaussian}.

\subsection{A finite precision requirement rather than a free real tape}
For $d=1$ the invariant domain is an interval inside the probability
simplex. A codebook can be constructed using interval intersections, and
its labels are ordinary finite integers. One evaluates the two rational
Bayes fractions, with the command likelihood evaluated to a specified
accuracy, then selects the nearest admissible representative. The
perturbation form of Theorem~\ref{thm:emulation} accounts for the incurred
state error. Since denominators are at least $3/4$, no division by a
vanishing report probability occurs.

For general fixed $d$, the same procedure uses a finite polytope
feasibility construction. It may be computationally expensive. Our
cardinality theorem does not state an optimal running time or workspace
bound. The program may store its predetermined codebook, as permitted by
the read-only model; it may not put acquired commands or reports in that
program. The current command must be discarded after updating the label,
and the terminal decoder may not reread it as free side information.

\subsection{Role of the accompanying diagnostics}
The source package includes exact-rational checks of Bayes normalization,
the odds inverse, the rank-one determinant identity, the refresh
contraction inequality at a finite family of inputs, and the scalar
quantization identity. It also checks a finite conditional-covariance
identity, a feedback path-law comparison, and specified deliberately
incorrect formulas. These computations are reproducibility aids for
signs, constants, and implemented examples.

The universal statements in this article are proved analytically above.
A finite test does not establish measurability, a global differential
bound, a minimax lower bound, or uniformity over all times. In particular,
no number of successful finite diagnostics certifies that every possible
hidden Markov model has the displayed memory law. Build receipts,
source hashes, and diagnostic output belong to the accompanying delivery
record rather than to the logical assumptions of any theorem.
